\subsection{对代数的应用}

设 A 是环，B 是 A-代数（不必结合或交换，也不必有单位元），S 是 A 的乘性子集。我们知道，可以在 \(S^{-1}A\)-模 \(S^{-1}B = B \otimes_{A} S^{-1}A\) 上定义一个典范的 \(S^{-1}A\)-代数结构，称为通过把标量环扩张到 \(S^{-1}A\) 而得到的（《代数学》第 III 章，§ 3），在其下乘积 \((x/s)(y/t)\) 等于 \((xy)/(st)\)。若 e 是 B 的单位元，则 e/1 是 \(S^{-1}B\) 的单位元，且若 B 是结合的（相应地，交换的），则 \(S^{-1}B\) 亦然。

设 A 是环，M 是 A-模；我们用 T(M)（相应地，∧(M)、S(M)）表示 M 的张量代数（相应地，外代数、对称代数）（《代数学》第 III 章）。已知对每个交换 A-代数 C，存在唯一的同构 \(j\)，从 \(\mathsf{T}(\mathbf{M}) \otimes_{\mathbf{A}} \mathbf{C}\) 到 \(\mathsf{T}(\mathbf{M} \otimes_{\mathbf{A}} \mathbf{C})\)（相应地，从 \(\bigwedge (\mathbf{M}) \oplus_{\mathbf{A}} \mathbf{C}\) 到 \(\bigwedge (\mathbf{M} \otimes_{\mathbf{A}} \mathbf{C})\)、从 \(\mathsf{S}(\mathbf{M}) \otimes_{\mathbf{A}} \mathbf{C}\) 到 \(\mathsf{S}(\mathbf{M} \otimes_{\mathbf{A}} \mathbf{C})\)），使得

\[
i (x \otimes 1) = x \otimes 1
\]

对所有 \(x \in \mathbf{M}\) 成立，这里 \(\mathbf{M}\) 典范地等同于 \(T(\mathbf{M})\)（相应地，\(\bigwedge (\mathbf{M}), S(\mathbf{M})\)）的子模（出处同前）。于是特别地可以看到，对每个乘性子集 \(\mathbf{R}\)（\(\mathbf{A}\) 的），有典范同构

\[
\begin{array}{c} \mathbf {R} ^ {- 1} \mathsf {T} (\mathbf {M}) \to \mathsf {T} (\mathbf {R} ^ {- 1} \mathbf {M}), \qquad \mathbf {R} ^ {- 1} \bigwedge (\mathbf {M}) \to \bigwedge (\mathbf {R} ^ {- 1} \mathbf {M}), \\ \mathbf {R} ^ {- 1} \mathsf {S} (\mathbf {M}) \to \mathsf {S} (\mathbf {R} ^ {- 1} \mathbf {M}) \end{array}
\]

它们在 \(R^{-1}M\) 上限制为恒同映射。

